\documentclass[11pt]{article}
\usepackage[a4paper,margin=19mm]{geometry}
\usepackage{fontspec}
\setmainfont{lmroman10-regular.otf}
\usepackage{amsmath,amssymb,amsthm,mathtools,xcolor,graphicx,hyperref}
\hypersetup{hidelinks}
\numberwithin{equation}{section}
\newtheorem{theorem}{Theorem}[section]
\newtheorem{lemma}[theorem]{Lemma}
\newtheorem{proposition}[theorem]{Proposition}
\newtheorem{corollary}[theorem]{Corollary}
\newtheorem{claim}[theorem]{Claim}
\newtheorem{fact}[theorem]{Fact}
\newtheorem{conjecture}[theorem]{Conjecture}
\theoremstyle{definition}
\newtheorem{definition}[theorem]{Definition}
\newtheorem{example}[theorem]{Example}
\newtheorem{remark}[theorem]{Remark}
\newtheorem{question}[theorem]{Question}
\theoremstyle{plain}
\newcommand{\C}{\mathbb C}
\newcommand{\R}{\mathbb R}
\newcommand{\Q}{\mathbb Q}
\newcommand{\Z}{\mathbb Z}
\newcommand{\N}{\mathbb N}
\newcommand{\eee}{\mathrm e}
\newcommand{\cclass}[1]{\textsf{#1}}
\newcommand{\cproblem}[1]{\textsc{#1}}
\newcommand{\ie}{i.\,e.}
\newcommand{\Ie}{I.\,e.}
\newcommand{\eg}{e.\,g.}
\newcommand{\Eg}{E.\,g.}
\newcommand{\phd}{Ph.\,D.}
\newcommand{\msc}{M.\,S.}
\newcommand{\bsc}{B.\,S.}
\newcommand{\expref}[2]{\hyperref[#2]{#1~\ref{#2}}}
\newcommand{\expeqref}[2]{\hyperref[#2]{#1~\eqref{#2}}}
\newcommand{\exprefsilent}[2]{\hyperref[#2]{#1}}
\providecommand{\etal}{et~al.}
\usepackage{eprint}
\usepackage{amssymb}
\usepackage{amsthm}
%

%

%
%
%
%
%
%
%
%
%
%

\newcommand{\srikanth}[1]{\todo[color=brown,size=\tiny]{S: #1}}
\newcommand{\ben}[1]{\todo[color=red!100!green!33, size=\footnotesize]{B: #1}}
%

\newcommand{\eps}{\varepsilon}
\newcommand{\mb}{\mathbf}
\newcommand{\tsfrac}[2]{{\textstyle\frac{#1}{#2}}}

%

%
%
%
%
%
%
%
%
%
%
%
%

%
\newcommand{\prob}[2]{\mathop{\mathrm{Pr}}_{#1}[#2]}
\newcommand{\avg}[2]{\mathop{\textbf{E}}_{#1}[#2]}
\newcommand{\poly}{\mathop{\mathrm{poly}}}
\newcommand{\tm}[1]{\textrm{#1}}
\newcommand{\bb}[1]{\mathbb{#1}}
\newcommand{\abs}[1]{\left|#1\right|}
\newcommand{\setcond}[2]{\left\{#1\: \middle|\: #2\right\}}
\newcommand{\errnote}[1]{{ \small #1}} 


%
\newcommand{\F}{\mathbb{F}}
\newcommand{\ip}[2]{\langle #1, #2 \rangle}


%
\newcommand{\Vbl}{\mathrm{Vbl}}
\newcommand{\Var}{\mathrm{Var}}
\newcommand{\Bf}{\mathbf{f}}
\newcommand{\mynorm}[1]{\lVert #1 \rVert}
\newcommand{\mc}[1]{\mathcal{#1}}
\newcommand{\dt}{\mathrm{D}}
\newcommand{\bs}{\mathrm{bs}}
\newcommand{\Maj}{\mathrm{Maj}}
\newcommand{\AC}{\cclass{AC}}
\newcommand{\NC}{\cclass{NC}}

\usepackage{enumerate}
\title{Approximate-majority gap between parity formulas and monotone circuits}
\author{Ben Rossman, Srikanth Srinivasan}
\date{doi:10.4086/toc.2019.v015a017; CC BY 3.0}
\begin{document}
\maketitle
\noindent\textit{Editorial note.} Full original author body/proof retained. Different original depth ranges and constants independent of d preserved. Randomized existence construction is not relabeled deterministic explicit. Upper/lower bounds grouped once.\par\medskip
\noindent\textit{Allowed standard prerequisites (editorial).}\begin{itemize}
\item Boolean circuits/formulas, probabilistic polynomials and elementary concentration. Original approximate-majority convention and all local approximation, bias amplification and circuit layers retained.. No collector reproof of cited standard polynomial-degree lower bounds; no deterministic-explicit upgrade of the original randomized circuit construction.
\end{itemize}\medskip
\section{Introduction}

The relative power of formulas versus circuits is one of the great mysteries in complexity \mbox{theory}. The central question in this area is whether $\NC^1$ (the class of languages decidable by polynomial-size Boolean formulas) is a proper subclass of P/poly (the class of languages decidable by polynomial-size Boolean circuits). Despite decades of efforts, this question remains wide open.\footnote{In this paper we focus on non-uniform complexity classes. The question of uniform-$\NC^1$ vs.\ $\cclass{P}$ is wide open as well.}
In the meantime, there has been progress on analogues of the $\NC^1$ vs.\ \cclass{P}/poly question in certain restricted settings. For instance, in the monotone basis (with AND and OR gates only), the power of polynomial-size formulas vs.\ circuits was separated by the classic lower bound of Karchmer and Wigderson \cite{karchmer1990monotone} (on the monotone formula size of st-Connectivity).

The bounded-depth setting is another natural venue for investigating
the question of %
formulas  %
vs.\ circuits. Consider the elementary fact that every depth-$d$ circuit of size $s$ is equivalent to a depth-$d$ formula of size at most $s^{d-1}$, where we measure \emph{size} by the number of gates. This observation is valid with respect to any basis (\ie, set of gate types). In particular, we may consider the $\AC^0$ basis (unbounded fan-in AND, OR, NOT gates) and the $\AC^0[\oplus]$ basis (unbounded fan-in MOD$_2$ gates in addition to AND, OR, NOT gates). With respect to either basis, there is a natural depth-$d$ analogue of the $\NC^1$ vs.\ \cclass{P}/poly question (where $d = d(n)$ is a parameter that may depend on $n$), namely whether every language decidable by polynomial-size depth-$d$ circuits is decidable by depth-$d$ formulas of size $n^{o(d)}$ (\ie, better than the trivial $n^{O(d)}$ upper bound).

It is reasonable to expect that this question could be resolved in the
%
sublogarithmic-depth  %
regime
%
$d(n) = o(\log n)$     %
given the powerful
%
lower-bound   %
techniques against $\AC^0$ circuits (H{\aa}stad's
Switching Lemma \cite{Hastad}) and $\AC^0[\oplus]$ circuits 
(the Polynomial Method of Razborov \cite{Razborov} and
Smolensky \cite{Smolensky87}). However, because the standard way
of applying these techniques does not distinguish between circuits
and formulas, it is not clear how to prove quantitatively stronger 
lower bounds on formula size vis-a-vis circuit size of a given function. 
Recent work of Rossman \cite{Rossman2015average} developed a new way of 
applying H{\aa}stad's Switching Lemma to $\AC^0$ formulas, in order to 
prove an $\exp(\Omega(dn^{1/(d-1)}))$ lower bound on the formula size 
of the Parity function for all $d \le O(\log n)$. Combined with the 
well-known $\exp(O(n^{1/(d-1)}))$ upper bound on the circuit size of 
Parity, this yields an asymptotically optimal separation in the 
power of depth-$d$ $\AC^0$ formulas vs.\ circuits for all 
$d(n) \le O({\log(n)}/{\log\log(n)})$, as well as a 
super-polynomial separation for all $\omega(1) \le d(n) \le o(\log n)$.

In the present paper, we carry out a similar development for formulas vs.\ circuits in the $\AC^0[\oplus]$ basis, obtaining both an asymptotically optimal separation for all $d(n) \le O({\log(n)}/{\log\log(n)})$ and a super-polynomial separation for all $\omega(1) \le d(n) \le o(\log n)$. Our target functions lie in the class of
\emph{Approximate Majorities}, here defined as Boolean 
functions $\{0,1\}^n \to \{0,1\}$ that  %
agree with  %
the Majority function on
a   %
$3/4$ fraction of
the  %
inputs. 
First, we show how to apply the Polynomial Method to obtain better parameters in the approximation of $\AC^0[\oplus]$ formulas by low-degree polynomials over $\F_2$. This leads to an $\exp(\Omega(dn^{1/2(d-1)}))$ lower bound on the $\AC^0[\oplus]$ formula size of all Approximate Majority functions. The other half of our formulas vs.\ circuits separation comes from an $\exp(O(n^{1/2(d-1)}))$ upper bound on the $\AC^0[\oplus]$ circuit size of some Approximate Majority function. In fact, this upper bound is realized by a randomized construction of monotone $\AC^0$ circuits
%
%
($\AC^0$ circuits without NOT or MOD$_2$ gates).   %



\begin{theorem}\label{thm:main}
Let $n \in \mathbb N$ be a growing parameter.
\begin{enumerate}
\item
For all $2 \le d \le O(\log n)$, the depth-$d$ $\AC^0[\oplus]$ formula size of any $n$-variable Approximate Majority function is $\exp(\Omega(dn^{1/2(d-1)}))$.
\item
  For all
  \[
    2 \le d \le \frac{1}{2}\cdot \frac{\log n}{\log\log n + O(1)},
  \]
  there exists an $n$-variable Approximate Majority function whose
  depth-$d$ monotone $\AC^0$ circuit size is $\exp(O(n^{1/2(d-1)}))$.
\end{enumerate}
\end{theorem}


(All asymptotic notation $O(\cdot)$ and $\Omega(\cdot)$ hides constants independent of $d$.) 
%
%
%
%
%
Denote by $\AC^0[\oplus]$-circuit$[d(n),s(n)]$
and $\AC^0[\oplus]$-formula$[d(n),s(n)]$
the class of languages decidable by $\AC^0[\oplus]$-circuits
and formulas, resp., of depth $d(n)$ and size $s(n)$.
%
For all functions $d(n)$, we have
%
the inclusions
\[
  \AC^0[\oplus]\textup{-formula}[d(n),\,\mathrm{poly}(n)]
  \subseteq
  \AC^0[\oplus]\textup{-circuit}[d(n),\,\mathrm{poly}(n)]
  \subseteq
  \AC^0[\oplus]\textup{-formula}[d(n),\,n^{O(d(n))}].  
\]
\expref{Theorem}{thm:main} implies that the above
%
inclusions       %
are tight for small $d(n)$.

\begin{corollary}\label{cor:separations}
Let $n \in \mathbb N$ be a growing parameter.
\begin{enumerate}[\normalfont(i)]
\item
For all functions $2 \le d(n) \le O\bigl(\frac{\log n}{\log \log n}\bigr)$, 
\[
  \AC^0[\oplus]\textup{-circuit}[d(n),\,\mathrm{poly}(n)]
  \nsubseteq
%
%
  \AC^0[\oplus]\textup{-formula}[d(n),\,n^{o(d(n))}].   %
\]
\item
For all functions $\omega(1) \le d(n) \le o(\log n)$, 
\[
  \AC^0[\oplus]\textup{-formula}[d(n),\,\mathrm{poly}(n)]
  \subsetneqq
  \AC^0[\oplus]\textup{-circuit}[d(n),\,\mathrm{poly}(n)].
\]
\end{enumerate}
\end{corollary}

\begin{proof}
  We prove separation (i) using a padding argument. Let $c > 0$ be a constant such that \expref{Theorem}{thm:main}(2) holds for all
  \[
    2 \le d \le \frac{1}{2}\cdot \frac{\log n}{\log\log n + c} + 1.
  \]
  And let $\{f_{n,d}\}$ be the family of Approximate Majority
  functions from \expref{Theorem}{thm:main}(2).  Now suppose
  \[
    2 \le d \le (1/2)^{c+1} \frac{\log n}{\log\log n} + 1.
  \]
  Letting $m = (\log n)^{2(d-1)}$, we have
\begin{align*}
  d - 1 &\le (1/2)^{c+1} \frac{\log n}{\log\log n}
  = (d-1) (1/2)^c \frac{m^{1/2(d-1)}}{\log m}.
\end{align*}
Dividing both sides by $d-1$ and taking logs, this rearranges to
\[
  d \le \frac{1}{2}\cdot \frac{\log m}{\log\log m + c} + 1.
\]
By \expref{Theorem}{thm:main}, $f_{m,d}$ has $\AC^0[\oplus]$ circuit size $\poly(n)$ and $\AC^0[\oplus]$ formula size $n^{\Omega(d)}$. Separation (i) now follows by padding $f_{m,d}$ to a function on $n$ variables (noting that $m < n$).

Separation (ii) follows from (i) when
\[
  d \le \frac{1}{2}\cdot \frac{\log n}{\log\log n + c} + 1.
\]
For larger depth, separation (ii) follows from
\expref{Theorem}{thm:main}, noting that the lower bound
$\exp(\Omega(dn^{1/2(d-1)}))$ is super-polynomial for
$d(n) = o(\log n)$.
\end{proof}

We point out that extending separation (ii) from depth $o(\log n)$ to depth $\log n$ is equivalent to separating complexity classes $\NC^1$ ($= \AC^0[\oplus]\textup{-formula}[\log n,\,\mathrm{poly}(n)]$) from $\AC^1$ ($= \AC^0[\oplus]\textup{-circuit}[\log n,\,\mathrm{poly}(n)]$). 


\subsection{Proof outline}

\paragraph{Improved polynomial approximation.} The lower bound for $\AC^0[\oplus]$ formulas follows the general template due to Razborov~\cite{Razborov} on proving lower bounds for $\AC^0[\oplus]$ circuits using low-degree polynomials over $\F_2$. Razborov showed that for any Boolean function $f:\{0,1\}^n\rightarrow \{0,1\}$ that has an $\AC^0[\oplus]$ circuit of size $s$ and depth $d$, there is a randomized polynomial $\mb P$ of degree $O(\log s)^{d-1}$ that computes $f$ correctly on each input with probability ${7}/{8}$ (we call such polynomials $1/8$-error \emph{probabilistic polynomials}). By showing that some explicit Boolean function $f$ (\eg, the Majority function or the $\mathrm{MOD}_q$ function for $q$ odd) on $n$ variables does not have such an approximation of degree less than $\Omega(\sqrt{n})$~\cite{Razborov,Smolensky87,Szegedy1989,  %
Smolensky93}, we get that any $\AC^0[\oplus]$ circuit of depth $d$ computing $f$ must have size $\exp(\Omega(n^{1/2(d-1)}))$.

In this paper, we improve the parameters of Razborov's polynomial approximation from above for $\AC^0[\oplus]$ \emph{formulas}. More precisely, for $\AC^0[\oplus]$ formulas of size $s$ and depth $d$, we are able to construct $1/8$-error probabilistic polynomials of degree $O(({1}/{d})\log s)^{d-1}$. (Since every depth-$d$ circuit of size $s$ is equivalent to a depth-$d$ formula of size at most $s^{d-1}$, this result implies Razborov's original theorem that $\AC^0[\oplus]$ circuits of size $s$ and depth $d$ have $1/8$-error probabilistic polynomials of degree $O(\log s)^{d-1}$.)

We illustrate the idea behind this improved polynomial approximation with the special case of a balanced formula (\ie, all gates have the same fan-in) of fan-in $t$ and depth $d$. Note that the size of the formula (number of gates) is $\Theta(t^{d-1})$ and hence it suffices in this case to show that it has a $1/8$-error probabilistic polynomial of degree $O(\log t)^{d-1}$. We construct the probabilistic polynomial inductively. Given a balanced formula $F$ of depth $d$ and fan-in $t$, let $F_1,\ldots,F_t$ be its subformulas of depth $d-1$. Inductively, each $F_i$ has a $1/8$-error probabilistic polynomial ${\mb P}_i$ of degree $O(\log t)^{d-2}$ and by a standard error-reduction~\cite{kopparty2012certifying}, it has a $(1/16t)$-error probabilistic polynomial of degree $O(\log t)^{d-1}$ (in particular, at any given input $x\in \{0,1\}^n$, the probability that \emph{there exists} an $i\in [t]$ such that ${\mb P}_i(x)\neq F_i(x)$ is at most $1/16$). Using Razborov's construction of a $1/16$-error probabilistic polynomial of degree $O(1)$ for the output gate of $F$ and composing this with the probabilistic polynomials ${\mb P}_i$, we get the result for balanced formulas. This idea can be extended to general (\ie, not necessarily balanced) formulas with a careful choice of the error parameter for each subformula $F_i$ to obtain the stronger polynomial approximation result. 

\paragraph{Improved formula lower bounds.} Combining the above approximation result with known lower bounds for polynomial 
approximation~\cite{Razborov,Smolensky87,Szegedy1989, %
Smolensky93}, we can already obtain stronger lower bounds for $\AC^0[\oplus]$ formulas than are known for $\AC^0[\oplus]$ circuits. For instance, it follows that any $\AC^0[\oplus]$ formula of depth $d$ computing the Majority function on $n$ variables must have size $\exp(\Omega(dn^{1/2(d-1)}))$ for all $d \le O(\log n)$, which is stronger than the corresponding circuit lower bound. Similarly stronger formula lower bounds also follow for the $\mathrm{MOD}_q$ function ($q$ odd). 

\paragraph{Separation between formulas and circuits.} However, the above improved lower bounds do not directly yield the claimed separation between $\AC^0[\oplus]$ formulas and circuits. This is because we do not have circuits computing (say) the Majority function of the required size. To be able to prove our result, we would need to show that the Majority function has $\AC^0[\oplus]$ circuits of depth $d$ and size $\exp(O(n^{1/2(d-1)}))$ (where the constant in the $O(\cdot)$ is independent of $d$). However, as far as we know, the strongest result in this direction~\cite{KPPY84majubd} only yields $\AC^0[\oplus]$ circuits of size greater than 
$\exp(\Omega(n^{1/(d-1)})),$\footnote{Indeed, this is inevitable with all constructions that we are aware of, since they are actually $\AC^0$ circuits and it is known by a result of H\r{a}stad~\cite{Hastad} that any $\AC^0$ circuit of depth $d$ for the Majority function must have size $\exp(\Omega(n^{1/(d-1)}))$.} which is superpolynomially larger than the upper bound.

To circumvent this issue, we change the hard functions to the class of \emph{Approximate Majorities}, which is the class of Boolean functions that agree with
the  %
%
%
%
Majority function on a $3/4$ fraction of inputs from $\{0,1\}^n$. (The choice of the constant $3/4$ is not crucial: it can be replaced by any constant in the interval $(1/2,1)$.)
%
While this has the downside that we no longer are dealing with an explicitly defined function, the advantage is that the polynomial approximation method of Razborov yields tight lower bounds for some functions from this class. 

Indeed, since the method of Razborov is based on polynomial approximations, it immediately follows that the same proof technique also yields the same lower bound for computing Approximate Majorities. Formally, any $\AC^0[\oplus]$ circuit of depth $d$ computing \emph{any} Approximate Majority must have size $\exp(\Omega(n^{1/2(d-1)}))$. On the upper bound side, it is known from the work of O'Donnell and Wimmer~\cite{odonnellwimmer} and Amano~\cite{Amano-appxmaj} that there \emph{exist} Approximate Majorities that can be computed by monotone $\AC^0$ \emph{formulas} of depth $d$ and size $\exp(O(dn^{1/2(d-1)}))$. (Note that the double exponent ${1}/{(2(d-1))}$ is now the same in the upper and lower bounds.)

In order to separate $\AC^0[\oplus]$ formulas and circuits, we show that both the upper and lower bounds from the previous paragraph are actually tight (up to the constants implicit in the exponents). Plugging in our stronger polynomial approximation for $\AC^0[\oplus]$ formulas, we obtain that any $\AC^0[\oplus]$ formula of depth $d$ computing any Approximate Majority must have size $\exp(\Omega(dn^{1/2(d-1)}))$. In particular, this implies that Amano's construction is tight (up to the universal constant in the exponent) even for $\AC^0[\oplus]$ formulas. 

Further, we also modify Amano's construction~\cite{Amano-appxmaj} to obtain better constant-depth \emph{circuits} for Approximate Majorities: we show that there exist Approximate Majorities that are computed by monotone $\AC^0$ circuits of depth $d$ of size $\exp(O(n^{1/2(d-1)}))$ (the constant in the $O(\cdot)$ is independent of $d$). 

\paragraph{Smaller circuits for Approximate Majority.} Our construction closely follows Amano's, which in turn is related to Valiant's probabilistic construction~\cite{Valiant-maj} of monotone formulas for the Majority function. However, we need to modify the construction in a suitable way that exploits the fact that we are constructing \emph{circuits}. This modification is in a similar spirit to a construction of Hoory, Magen and Pitassi~\cite{hoory2006monotone} who modify Valiant's construction to obtain smaller monotone circuits (of depth $\Theta(\log n)$) for computing the Majority function exactly.

At a high level, the difference between Amano's construction and ours is as
follows. Amano constructs random formulas $F_i$ of each depth $i\leq d$ as
follows. The formula $F_1$ is the AND of $a_1$ independent and randomly
chosen variables.
%
%
%
%
For even $i>1$, $F_i$ is the OR  of $a_i$ independent and random copies of $F_{i-1}$, while for odd $i > 1$, $F_i$ is the AND  of $a_i$ independent and random copies of $F_{i-1}$.
%
For suitable values
of $a_1,\ldots,a_d\in \mathbb{N}$, the random formula $F_d$ computes an
Approximate Majority with high probability. In our construction, we build
a depth $i$ circuit $C_i$ for each $i\leq d$ in a similar way, except that
each $C_i$ now has $M$ different outputs. Given such a $C_{i-1}$, we
construct $C_i$ by taking $M$ independent randomly chosen subsets
$T_1,\ldots,T_M$ of $a_i$ many outputs of $C_{i-1}$ and adding gates
that compute either the OR or AND (depending on whether $i$ is even
or odd) of the gates in $T_i$. Any of the $M$ final gates of $C_d$ now
serves as the output gate. By an analysis similar to Amano's
(see also~\cite{hoory2006monotone}) we can show that this computes an
Approximate Majority with high probability, which finishes
the proof.\footnote{This is a slightly imprecise description of the
construction as the final two levels of the circuit are actually defined
somewhat differently.}



\section{Preliminaries}

Throughout, $n$ will be a growing parameter. We will consider Boolean functions on $n$ variables, \ie, functions of the form $f:\{0,1\}^n\rightarrow \{0,1\}$. We will sometimes identify $\{0,1\}$ with the field $\F_2$ in the natural way and consider functions $f:\F_2^n \rightarrow \F_2$ instead.

%
%
%
%
%
%
%
%

Given a Boolean vector $y\in \{0,1\}^n$, we use $|y|_0$ to denote the number of $0$s in $y$ and $|y|_1$ to denote the number of $1$s in $y$.
%
%


The Majority function on $n$ variables, denoted $\mathrm{MAJ}_n$ is the
Boolean function that maps inputs $x\in \{0,1\}^n$ to $1$ if and
only if $|x|_1 > n/2$. 

\begin{definition}
\label{defn:approx-maj}
%
%
An \emph{$(\varepsilon,n)$-Approximate Majority} is a function $f:\{0,1\}^n\rightarrow \{0,1\}$ such that
\[
  \prob{x \in \{0,1\}^n}{f(x) \neq \Maj_n(x)} \leq \varepsilon.
\]
\end{definition}

As far as we know, the study of this class of functions was initiated by O'Donnell and Wimmer~\cite{odonnellwimmer}. See also~\cite{Amano-appxmaj, BlaisTan}.

We refer the reader to~\cite{AroraBarak, Jukna} for standard definitions of
Boolean circuits and formulas. We use
%
the terms  %
%
%
\emph{$\AC^0$ circuits} and \emph{$\AC^0$ formulas} to describe
%
%
and formulas %
of constant depth,   %
made up of
$\mathrm{AND},\ \mathrm{OR}$ and $\mathrm{NOT}$ gates. 
Similarly, %
\emph{$\AC^0[\oplus]$ circuits}
%
and %
%
\emph{$\AC^0[\oplus]$ formulas}
will be circuits
%
%
and formulas %
of constant depth,   %
made up of $\mathrm{AND},\ \mathrm{OR},\ \mathrm{MOD}_2$ 
and $\mathrm{NOT}$ gates.
%
%
\begin{remark}
%
%
%
%
In standard terminology, $\AC^0$ and $\AC^0[\oplus]$ circuits refer to circuits of \emph{constant} depth and \emph{polynomial} size (made up of the relevant sets of gates in each case). However, our results hold for more general settings of these parameters. To phrase these results, it will be convenient to abuse notation and refer to $\AC^0$ and $\AC^0[\oplus]$ circuits of size $s(n)$ and depth $d(n)$ where $s(n)$ is possibly superpolynomial and $d(n)$ possibly superconstant, which of course can be defined in the obvious way.
\end{remark}
%
%
%
%
%
By the \emph{size of a circuit} we mean the number of \emph{gates}
in the circuit.  (Note that this, too, is a deviation from standard
terminology that refers to the number of wires as the ``size'' of
%
%
the circuit. However, since the number of wires is polynomially related to the number of gates, \expref{Theorem}{thm:main} and \expref{Corollary}{cor:separations} also hold when the size denotes the number of wires.)
%
By the \emph{size of a formula} we mean the number of its \emph{leaves}
%
which is, %
up to a constant factor, equal to
%
the number of all nodes (including leaves)
in the formula.\footnote{We assume here without loss of generality
that the formula does not contain a gate of fan-in $1$ feeding into
%
another gate.}  %

%
\section{Lower bound}    %
\label{sec:lbd}

In this section, we show that any $\AC^0[\oplus]$ formulas of depth $d$
computing a $(1/4,n)$-Approximate Majority must have size at least
$\exp(\Omega(dn^{1/2(d-1)}))$ for all $d \le O(\log n)$. 

We work over the field $\F_2$ and identify it with $\{0,1\}$ in the
natural way. The following concepts are standard in circuit complexity
(see, \eg, Beigel's survey~\cite{Beigelpolysurvey}).

\begin{definition}
Fix any $\varepsilon\in [0,1]$. A polynomial $P\in \F_2[X_1,\ldots,X_n]$ is
said to be an \emph{$\varepsilon$-approximating polynomial} for a Boolean
function $f:\{0,1\}^n\rightarrow \{0,1\}$ if 
\[
\prob{x\in \{0,1\}^n}{f(x) = P(x)}\geq 1-\varepsilon.
\]
\end{definition}

%
The following lemma, motivated by the ideas of 
Smolensky~\cite{Smolensky87} (1987), first appeared explicitly in
Szegedy's Ph. D. Thesis~\cite[Theorem 2.4.7]{Szegedy1989} (1989)
and in Smolensky's paper~\cite[Theorem 5]{Smolensky93} (1993).

%
%
 %

\begin{lemma}[Smolensky, Szegedy]
\label{lem:razb-maj}
Let $\varepsilon \in (0,{1}/{2})$ be any fixed constant. Any $(1/2-\varepsilon)$-approximating polynomial for the Majority function on $n$ variables must have degree $\Omega(\sqrt{n})$. 
\end{lemma}

\begin{corollary}
\label{cor:approx-maj-deg}
Let $f$ be any $(1/4,n)$-Approximate Majority and $\varepsilon\in (0,1/4)$ an arbitrary constant. Then any $({1}/{4}-\varepsilon)$-approximating polynomial for $f$ must have degree $\Omega(\sqrt{n})$.
\end{corollary}

\begin{proof}
The proof is immediate from \expref{Lemma}{lem:razb-maj} and the triangle inequality.
\end{proof}

\begin{definition}
An \emph{$\eps$-error probabilistic polynomial of degree $D$} for a Boolean function $f : \{0,1\}^n \to \{0,1\}$ is a random variable $\mb P$ taking values from polynomials in $\F_2[X_1,\ldots,X_n]$ of degree at most $D$ such that for all $x \in \{0,1\}^n$, we have $\Pr[\ f(x) = \mb P(x)\ ] \ge 1 - \eps$.
\end{definition}

\begin{definition}
Let $D_\eps(f)$ be the minimum degree of an $\eps$-error probabilistic polynomial for $f$.
\end{definition}

We will make use of the following two lemmas concerning $D_\eps(\cdot)$.

\begin{lemma}[Razborov~\cite{Razborov}]
\label{lem:razb-OR}
Let $\mathrm{OR}_n$ and $\mathrm{AND}_n$ be the OR and AND functions
%
in    %
$n$ variables, %
respectively.
%
Then $D_\varepsilon(\mathrm{OR}_n), D_\varepsilon(\mathrm{AND}_n)\leq \lceil \log(1/\varepsilon)\rceil$.
\end{lemma}

\begin{lemma}[Kopparty and Srinivasan~\cite{kopparty2012certifying}]\label{la:eps}
There is an absolute constant $c_1$ such that for any $\varepsilon \in (0,1)$, $D_\eps(f) \le c_1\cdot\lceil \log(1/\eps)\rceil \cdot D_{1/8}(f)$ for all Boolean functions $f$.
\end{lemma}

We now state our main result, which shows that every $\AC^0[\oplus]$ formula of size $s$ and depth $d + 1$ admits a $1/8$-error approximating polynomial of degree $O((1/d)\log s)^d$.

\begin{theorem}
  There is an absolute constant $c_2$ such that, if $f$ is computed by an $\AC^0[\oplus]$ formula $F$ of size $s$ and depth $d+1$, then
  \[
    D_{1/8}(f) \le 3\left(c_2\left(\frac{1}{d}\log(s)+1\right)\right)^d.
  \]
\end{theorem}

\begin{proof}
The proof is an induction on the depth $d$ of the formula. 

The base case $d=0$ corresponds to the case when the formula is a single AND, OR or MOD$_2$  gate and we need to show that $D_{1/8}(f)\leq 3$. In the case that the formula is an AND or OR gate, this follows from \expref{Lemma}{lem:razb-OR}. If the formula is a MOD$_2$  gate, this follows from the fact that the MOD$_2$  function is exactly a polynomial of degree $1$. 

Let $d\ge 1$. We assume that the formula $F$ is the AND/OR/MOD$_2$  of
%
subformulas   %
$F_1,\ldots,F_m$ computing $f_1,\dots,f_m$ where $F_i$ has size $s_i$ and depth $d+1$. So $F$ has size $s = s_1+\dots+s_m$ and depth $d+2$. Assume that
\[
  D_{1/8}(f_i) \le 3\left(c_2\left(\frac{1}{d}\log(s_i)+1\right)\right)^d
\]
for all $i$. We must show that
\[
  D_{1/8}(f) \le 3\left(c_2\left(\frac{1}{d+1}\log(s)+1\right)\right)^{d+1}.
\]

By \expref{Lemma}{la:eps}, each $f_i$ has an $s_i/(16s)$-error probabilistic polynomial $\mb P_i$ of degree $c_1\cdot\lceil\log(16s/s_i)\rceil\cdot D_{1/8}(f_i)$, which is at most
\[
  3c_1 \cdot 5(\log(s/s_i)+1)\cdot \left(c_2\left(\frac{1}{d}\log(s_i)+1\right)\right)^d.
\]
Then $(\mb P_1,\dots,\mb P_m)$ jointly computes $(f_1,\dots,f_m)$ with error $1/16$ ($= \sum_{i=1}^m (s_i/(16 s))$). 

By a reasoning identical to the base case, it follows that there exists a $1/16$-error probabilistic polynomial $\mb Q$ of degree $4$ for the output gate of the formula. 

Then $\mb Q(\mb P_1,\dots,\mb P_m)$ is a $1/8$-error probabilistic polynomial for $f$ of degree
\[
  60c_1\cdot\max_i\left\{ (\log(s/s_i)+1)\right\} \cdot \left(c_2\left(\frac{1}{d}\log(s_i)+1\right)\right)^d.
\]

So long as $c_2 \ge 20c_1$, it suffices to show that for all $i$,
\[
  (\log(s/s_i)+1) \cdot \left(\frac{1}{d}\log(s_i)+1\right)^d
  \le
  \left(\frac{1}{d+1}\log(s)+1\right)^{d+1}.
\]
Consider any $i$ and let $a,b \ge 0$ such that $s_i = 2^a$ and $s = 2^{a+b}$. We must show
\[
  (b+1)\left(\frac{a\vphantom{b}}{d}+1\right)^d
  \le
  \left(\frac{a+b}{d+1}+1\right)^{d+1}.
\]
For fixed $a \ge 0$, as a polynomial in $b$, the function
\[
  p_{a,d}(b) := \left(\frac{a+b}{d+1}+1\right)^{d+1} - (b+1)\left(\frac{a\vphantom{b}}{d}+1\right)^d
\]
is nonnegative over $b \ge 0$ with a unique root at $b = a/d$. This follows from
\[
  \frac{\partial}{\partial b} p_{a,d}(b) = 
  \left(\frac{a+b}{d+1}+1\right)^d - \left(\frac{a\vphantom{b}}{d}+1\right)^d,
\]
which is zero iff $b = a/d$; this value is a minimum of $p_{a,d}$ with $p_{a,d}(a/d) = 0$.
\end{proof}

\begin{corollary}
\label{cor:formula-lbd}
Let $n\in \mathbb{N}$ be a growing parameter and $d = d(n)$ (possibly a constant). 
Let $f$ be any $(1/4,n)$-Approximate Majority. Then any $\AC^0[\oplus]$ formula of depth $d$ computing $f$ must have size $\exp(\Omega(dn^{1/2(d-1)}))$ for all $d \le O(\log n)$, where the asymptotic notations $O(\cdot)$ and $\Omega(\cdot)$ hide absolute constants (independent of $d$ and $n$). 
\end{corollary}

\begin{proof}
Say that $F$ is an $\AC^0[\oplus]$ formula of depth $d$ and size $s$ computing $f$. Then, by \expref{Lemma}{la:eps}, we see that $F$ has a $1/8$-error probabilistic polynomial $\mb P$ of degree $D\leq O(({1}/{d})\log s+1)^{d-1}$. In particular, by an averaging argument, there is some fixed polynomial $P\in \F_2[X_1,\ldots,X_n]$ of degree at most $D$ such that $P$ is a $1/8$-error approximating polynomial for $f$. 

\expref{Corollary}{cor:approx-maj-deg} implies that the degree of $P$ must be $\Omega(\sqrt{n})$. Hence, we obtain $O(({1}/{d})\log s+1)^{d-1} \ge \Omega(\sqrt{n})$. It follows that
$s \ge \exp(\Omega(dn^{1/2(d-1)}) - O(d))$.
Observe that $\Omega(dn^{1/2(d-1)})$ dominates $O(d)$ so long as $d \le \eps \log n$ for some absolute constant $\eps > 0$ (depending on the constants in $\Omega(\cdot)$ and $O(\cdot)$). Hence, we get the claimed lower bound  $s \ge \exp(\Omega(dn^{1/2(d-1)}))$
for all $d \le \eps\log n$.
\end{proof}


%
\section{Upper bound}      %

In this section, we show that for any constant $\varepsilon$, there are $(\varepsilon,n)$-Approximate Majorities that can be computed by depth-$d$ $\AC^0$ circuits of size $\exp(O(n^{1/2(d-1)}))$.

\subsection{High-level idea}

We start with an informal description of our construction, which closely follows the constructions of O'Donnell and Wimmer~\cite{odonnellwimmer} and Amano~\cite{Amano-appxmaj} (see also~\cite{Valiant-maj,hoory2006monotone}).

The first basic observation, from \cite{odonnellwimmer}, is that to compute an $(\varepsilon,n)$-Approximate majority, it suffices to compute the Majority function on most (say at least a $1-(\varepsilon/4)$ fraction of) inputs of relative Hamming weight bounded away from $1/2$; say those of weight either at most ${1}/{2}-(\varepsilon /4\sqrt{n})$ or at least ${1}/{2}+(\varepsilon/4\sqrt{n})$. So we assume that the inputs come from one of these sets. Let us call the sets of high-weight inputs and low-weight inputs $Y$ (for ``Yes'') and 
$N$ (for ``No''),  %
respectively. 
%

The problem of distinguishing inputs from $Y$ and $N$ is easier than computing Majority exactly since there is a gap between the weights of the yes and no instances. In particular, given any inputs $a\in Y$ and $b\in N$, we have that the multiplicative gap between the weights of $a$ and $b$ is $(1+\gamma_0),$ where $\gamma_0 = \Theta(\varepsilon/\sqrt{n})$.

The idea is to iteratively increase this gap at each level of the circuit, so that after $(d-1)$ levels, the multiplicative gap is a large constant, at which point a simple OR or AND determines the answer on most inputs. 

The basic observation that allows us to do this (formalized in \expref{Lemma}{lem:technicalORAND} below) is the following. Say we have an input $x\in Y'\cup N'$ where $Y'$ is the set of inputs of relative weight at least $p(1+\gamma)$ and $N'$ is the set of inputs of relative weight at most $p(1-\gamma)$ (for some known $p,\gamma$). Consider the depth-$1$ circuit $C'$ that chooses $M$
%
random subsets of the input bits, %
%
%
of size $t$ each, 
and %
computes   %
the OR of the variables in each set ($C'$ is thus an $M$-output circuit). If the input $x$ has (relative) weight at least $p(1+\gamma)$, then the expected fraction of $0$s in the output---let us call this the \emph{$0$-weight} of the output---is at most
\[
(1-p(1+\gamma))^t\approx \exp(-pt -p\gamma t)
\]
while if $x$ has weight at most $p(1-\gamma),$ then the output has $0$-weight at least $\exp(-pt + p\gamma t)$.  By a simple application of the Chernoff bound and averaging, we can easily ensure that we have a fixed circuit $C'$ that has (roughly) this behaviour at every input $x$ of interest. That is, for inputs $a\in Y'$ and $b\in N'$, the $0$-weights of $C'(a)$ and $C'(b)$ have $0$-weights that are either at most $\exp(-pt)\cdot \exp(-\gamma pt)$ or at least $\exp(-pt)\cdot \exp(\gamma pt)$. In particular, they are multiplicatively separated by $1+\gamma'$ where $\gamma'\approx (pt)\cdot \gamma$,\footnote{We use the approximation that $\exp(z) \approx (1+z)$ when $z$ is small.} which is much
%
greater   %
than $\gamma$ if
%
$pt$ is much greater than $1$.   %
In an analogous way, we can construct a depth-$1$ circuit (made up of ANDs this time) that amplifies a multiplicative gap of $(1+\gamma)$ in the $0$-weights to a multiplicative gap of $(1+\gamma')$ in the weights. 

Given this fact, our way forward is clear. We stack these depth-$1$ circuits one on top of the other, at each stage amplifying the multiplicative gap from $1+\gamma$ (for some $\gamma$) to $1+r\gamma$.\footnote{One could also amplify the gaps by varying amounts at different levels of the circuit. This turns out to be %
%
%
below optimal 
in terms of the size of the circuit constructed.} Since we want the gap after $d-1$ levels to be a constant, we take $r \approx (1/\gamma_0)^{1/(d-1)} \approx n^{1/2(d-1)}$.

This in turns determines the size of our constructed circuit. Recall that the depth-$1$ circuit describes above amplifies the gap between the weights from $1+\gamma$ to $1+(pt)\cdot \gamma$, which we would like to be $1+r\gamma$. At the first level of the circuit, we can take $t\approx r$ as $p = 1/2 = \Theta(1)$. At all subsequent levels, however, we have $p\approx \exp(-r)$ and hence $t \approx r/p$ is $\exp(O(r))$. Hence, we obtain a circuit of size $\exp(O(r)) = \exp(O(n^{1/2(d-1)}))$.


\subsection{Formal proof}

Let $\varepsilon_0\in (0,1)$ be a small enough constant so that the following inequalities hold for any $\beta \leq \varepsilon_0$.
\begin{itemize}
\item $\exp(-\beta) \leq 1-\beta\exp(-\beta)$, and
\item $1-\beta \geq \exp(-\beta-\beta^2)\geq \exp(-2\beta)$.
\end{itemize}
(It suffices to take $\varepsilon_0 = 1/2$.) 

We need the following technical lemma.

\begin{lemma}
\label{lem:technicalORAND}
Let $A,s$ be positive reals, $M,n\in\mathbb{N}$,  and $\gamma\in ({1}/{n},{1}/{10})$ be such that $\eee^A \geq n^3$, $n \geq {1}/{\varepsilon_0}$, and $1\leq s\leq n$. Define
\[
  I_0(\gamma) := \{y\in \{0,1\}^M\ |\ |y|_1\leq M\eee^{-A}(1-\gamma)\}\quad\text{and}\quad I_1(\gamma) := \{y\in \{0,1\}^M\ |\ |y|_1\geq M\eee^{-A}(1+\gamma)\}.
\]
If we choose $S\subseteq [M]$ of size
$t := \lceil \eee^{A}\cdot s\rceil$ by picking $t$ random elements
from $M$ with replacement, then
\begin{align*}
x\in I_0(\gamma) &\Rightarrow \Pr_{S}\left[\bigvee_{j\in S}x_j = 0\right] \geq \exp(-s)\cdot\exp(s\gamma/2), \\
x\in I_1(\gamma) &\Rightarrow \Pr_{S}\left[\bigvee_{j\in S}x_j = 0\right] \leq \exp(-s)\cdot \exp(-s\gamma).
\end{align*}
Further, if $s\gamma \leq \varepsilon_0$, then the above probabilities can be 
%
bounded from below and from above  %
by $\exp(-s)\cdot (1+s\gamma\exp(-s\gamma))$ and
$\exp(-s)\cdot (1-s\gamma\exp(-s\gamma))$, %
respectively.  %

A similar statement can be obtained above for the sets
\[
  J_1(\gamma) := \{y\in \{0,1\}^M\ |\ |y|_0\leq M\eee^{-A}(1-\gamma)\}\quad\text{and}\quad J_0(\gamma) := \{y\in \{0,1\}^M\ |\ |y|_0\geq M\eee^{-A}(1+\gamma)\},
\]
with the event ``$\bigvee_{j\in S}x_j = 0$'' being replaced by the
event ``$\bigwedge_{j\in S}x_j = 1$.''
\end{lemma}

\begin{proof}
We give the proof only for $I_0(\gamma)$ and $I_1(\gamma)$. The proof for $J_0(\gamma)$ and $J_1(\gamma)$ is similar.

Consider first the case that $x\in I_1(\gamma)$. In this case, we have the following computation. 
\begin{align}
\prob{S}{\bigvee_{j\in S} x_j = 0} &\leq \left(1-\frac{1+\gamma}{\eee^A}\right)^{\eee^A\cdot s}
%
%
\leq \exp(-(1+\gamma)\cdot s) \leq \exp(-s)\cdot \exp(-s\gamma).\label{eq:tech2}
\end{align}
The above implies the first upper bound on $\prob{S}{\bigvee_{j\in S} x_j = 0}$ from the lemma statement. When $s\gamma\leq \varepsilon_0$, we further have $\exp(-s\gamma) \leq 1-s\gamma \exp(-s\gamma)$, which implies the second upper bound. This proves the lemma when $x\in I_1(\gamma)$.

Now consider the case that $x\in I_0(\gamma)$. We have
\begin{align}
  \prob{S}{\bigvee_{j\in S}x_j = 0} &\geq \left(1-\frac{1-\gamma}{\eee^A}\right)^{\eee^A\cdot s + 1}\notag\\
                                    &\geq \exp\left(\left(-\frac{1-\gamma}{\eee^A} - \frac{1}{\eee^{2A}}\right)\cdot (\eee^A\cdot s+1)\right)\notag\\
                                    &= \exp\left(-s+s\gamma -\frac{s}{\eee^A}-\frac{1-\gamma}{\eee^A} - \frac{1}{\eee^{2A}}\right)\notag\\
                                    &\geq \exp\left(-s+s\gamma -\frac{2s}{\eee^A}\right)\notag\\
                                    &= \exp(-s)\cdot \exp\left(s\gamma\left(1-\frac{2\eee^{-A}}{\gamma}\right)\right)\label{eq:tech1}
\end{align}
where for the second inequality we have used the fact that since $\eee^{-A} \leq {1}/{n^3}\leq \varepsilon_0$, we have
\[
  1- \frac{1-\gamma}{\eee^A}\geq \exp\left(-\frac{1-\gamma}{\eee^A} - \frac{1}{\eee^{2A}}\right).
\]
Since $\eee^{-A}\leq {1}/{n^3}\leq {1}/{(4n)}\leq \gamma/4$,
we can
%
%
bound from below
%
the right hand side of \expeqref{inequality}{eq:tech1} by $\exp(-s)\cdot \exp(s\gamma/2)$. Also, note that 
\begin{align*}
1-\frac{2\eee^{-A}}{\gamma} &\geq 1-\frac{2/n^3}{1/n}= 1-\frac{2}{n^2}
\geq \exp(-1/n) \geq \exp(-\gamma)\geq \exp(-s\gamma).
\end{align*}
This implies that the righthand side of \expeqref{inequality}{eq:tech1} can also be bounded from below by
\[
  \exp(-s)\exp(s\gamma\exp(-s\gamma)) \geq \exp(-s)\cdot(1+s\gamma\exp(-s\gamma)),
\]
which implies the claim about $\prob{S}{\bigvee_{j\in S} x_j = 0}$
assuming that $x\in I_0(\gamma)$.
\end{proof}


We now prove the main result of this section. 

\begin{theorem}
\label{thm:amano-ckt}
For any growing parameter $n\in \mathbb{N}$ and
\[
  2 \le d \le \frac{1}{2}\cdot \frac{\log n}{\log\log n + O(1)}
\]
and $\eps > 0$, there is an $(\varepsilon,n)$-Approximate Majority
$f_n$ computable by a monotone $\AC^0$ circuit with at most
$\exp(O(n^{1/2(d-1)}\log(1/\varepsilon)/\varepsilon))$ many gates,
where both $O(\cdot)$'s hide absolute constants (independent of
$d,\varepsilon$).
\end{theorem}

\begin{proof}
  We assume throughout that $\varepsilon$ is a small enough constant and that $n$ is large enough for various inequalities to hold. We will actually construct a monotone circuit of depth $d$ and size
  \[
    \exp\left(O\left(n^{1/2(d-1)}\frac{\log(1/\varepsilon)}{\varepsilon}\right)\right)
  \]
  computing a $(4\varepsilon,n)$-Approximate Majority, which also implies the theorem. 

  Fix parameter $A= n^{1/2(d-1)} - \alpha$ where $\alpha \in [0,\ln 2]$ is chosen so that $A = A_1\cdot \ln 2$ for some positive integer $A_1$. We assume that $A\geq 10\log n$ (which holds as long as
  \[
    2d \le \frac{\log n}{\log\log n + c}
  \]
  for a large enough absolute constant $c > 0$) and that
  $\varepsilon \leq \varepsilon_0$. Let $M = \lceil \eee^{10A}\rceil$.


Define a sequence of real numbers $\gamma_0,\gamma_1,\ldots,\gamma_{d-2}$ as follows.
\begin{align*}
\gamma_0 &= \frac{\varepsilon}{\sqrt{n}},\\
\gamma_i &= A\gamma_{i-1}\exp(-2A\gamma_{i-1}), \text{ for each $i\in [d-2]$.}
\end{align*}

It is clear that $\gamma_i \leq A^i \gamma_0$ for each $i\in [d-2]$. As a result we also obtain 
\begin{align}
\gamma_i &= A^i \gamma_0\exp(-2A(\gamma_0+\gamma_1+\cdots+\gamma_{i-1}))\notag\\
&\geq A^i \gamma_0\exp(-2\gamma_0A(1+A+A^2+\cdots+A^{i-1})) \geq A^i\gamma_0\exp(-3A^i\gamma_0).\label{eq:gamma-i}
\end{align}

Let
\begin{align*}
Y_\varepsilon &= \left\{x\in \{0,1\}^n\ \middle|\ |x|_1\geq \left(\frac{1}{2}+\frac{\varepsilon}{\sqrt{n}}\right)n\right\},\\
N_\varepsilon &= \left\{x\in \{0,1\}^n\ \middle|\ |x|_1\leq \left(\frac{1}{2}-\frac{\varepsilon}{\sqrt{n}}\right)n\right\}.
\end{align*}

The idea is to define a sequence of circuits $C_1,C_2,\ldots,C_{d-2}$ with $n$ inputs and $M$ outputs such that $C_i$ has depth $i$ and $iM$ many (non-input) gates. Further, for odd $i$, 
\begin{equation}
x\in N_\varepsilon \Rightarrow C_i(x)\in I_0(\gamma_i),\quad\text{ and }\quad
x\in Y_\varepsilon \Rightarrow C_i(x)\in I_1(\gamma_i),\label{eq:oddi}
\end{equation}
and similarly for even $i$,
\begin{equation}
x\in N_\varepsilon \Rightarrow C_i(x)\in J_0(\gamma_i),\quad\text{ and }\quad
x\in Y_\varepsilon \Rightarrow C_i(x)\in J_1(\gamma_i).\label{eq:eveni}
\end{equation}

After this is done, we will add on
the %
top a depth-$2$ circuit that will reject most inputs from $I_0(\gamma_{d-2})$
or $J_0(\gamma_{d-2})$---depending on whether
%
%
%
$(d-2)$ is odd or even---and accept most inputs from $I_1(\gamma_{d-2})$ or $J_1(\gamma_{d-2})$. 

We begin with the construction of $C_1,\ldots,C_{d-2}$ which is done by induction. 

\paragraph{Construction of $C_1$.} The base case of the induction is the construction of $C_1$, which is done as follows. We choose $M$ i.i.d.\ random subsets $T_1,\ldots,T_M\subseteq [n]$ in the following way: for each $i\in [M]$, we sample $A_1$ random elements of $[n]$ with replacement. Let $b_i^x = \bigwedge_{j\in T_i}x_j$.

If $x\in N_\varepsilon$, then the probability that $b_i^x = 1$ is given by 
\begin{align*}
\prob{}{b_i^x = 1} &\le \left(\frac{1}{2}- \gamma_0\right)^{A_1}
= \frac{1}{2^{A_1}} \left(1-2\gamma_0\right)^{A_1} \leq \frac{1}{\eee^A}\left(1 -\gamma_0A_1 \right) \leq \frac{1}{\eee^A}\left(1 -\gamma_0A\right),
\end{align*}
where the second-last inequality follows from the fact that
\[
  (1-z)^{A_1} \leq \left(1-zA_1 + \frac{A_1^2z^2}{2}\right)
\]
for $z\in [0,1]$.

Let $\delta = 1/n^3$. Note in particular that $2\delta/\gamma_0A \leq \varepsilon_0$ for large enough $n$.

By a Chernoff bound, the probability that
\[
  \frac{1}{M}\sum_i b_i^x\geq \frac{1}{\eee^A}(1-\gamma_0A)(1+\delta)
\]
is bounded by
$\exp(-\Omega(\delta^2M/\eee^A))\leq \exp(-\Omega(\eee^{9A}/n^6))\leq
\exp(-n)$, since $\eee^A \geq n^{10}$. Thus, with probability at least
$1-\exp(-n)$, we have
\begin{align}
\frac{\sum_i b_i^x}{M}&\leq \frac{1}{\eee^A}(1-\gamma_0A)(1+\delta)\notag\\
&\leq \frac{1}{\eee^A}(1-\gamma_0A + \delta) = \frac{1}{\eee^A}\left(1-\gamma_0A\left(1-\frac{\delta}{\gamma_0A}\right)\right)\notag\\
&\leq \frac{1}{\eee^A}\left(1-\gamma_0A\exp\left(-\frac{2\delta}{\gamma_0A}\right)\right)\leq \frac{1}{\eee^A}(1-\gamma_0A\exp(-\gamma_0A))\notag\\
&\leq \frac{1}{\eee^A}(1-\gamma_1).\label{eq:C1.1}
\end{align}
%
Above, we have used the fact that
\[
  \left(1-\frac{\delta}{\gamma_0A}\right)\geq \exp\left(\frac{-2\delta}{\gamma_0A}\right)
\]
since $\delta/\gamma_0A \leq \varepsilon_0$.

If $x\in Y_\varepsilon$, then the probability that $b_i^x = 1$ is given by 
\begin{align*}
\prob{}{b_i^x = 1} &\geq \left(\frac{1}{2}+ \gamma_0\right)^{A_1}
%
%
\geq \frac{1}{2^{A_1}} \left(1+2\gamma_0\right)^{A_1} 
\geq \frac{1}{\eee^A}(1+\gamma_0{A_1})
%
%
\geq \frac{1}{\eee^A}(1+\gamma_0A).
\end{align*}

As above, we can argue that the probability that
\[
  \frac{1}{M}\sum_i b_i^x \leq \frac{1}{\eee^A}(1+\gamma_0A)(1-\delta)
\]
is at most $\exp(-n)$. Thus, with probability $1-\exp(-n)$,
\begin{align}
\frac{\sum_i b_i^x}{M}&\geq \frac{1}{\eee^A}(1+\gamma_0A)(1-\delta)\notag\\
&\geq \frac{1}{\eee^A}(1+\gamma_0A -2\delta) = \frac{1}{\eee^A}\left(1+\gamma_0A\left(1-\frac{2\delta}{\gamma_0A}\right)\right)\notag\\
&\geq \frac{1}{\eee^A}\left(1+\gamma_0A\exp\left(-\frac{4\delta}{\gamma_0A}\right)\right)\geq \frac{1}{\eee^A}(1+\gamma_0A\exp(-\gamma_0A))\notag\\
&\geq \frac{1}{\eee^A}(1+\gamma_1).\label{eq:C1.2}
\end{align}
Thus, by a union bound over $x$, we can fix a choice of $T_1,\ldots,T_M$ so that \expeqref{inequality}{eq:C1.1} holds for all $x\in N_\varepsilon$ and \expeqref{inequality}{eq:C1.2} holds for all $x\in Y_\varepsilon$. Hence, \expeqref{condition}{eq:oddi} holds for $i=1$ as required. This concludes the construction of $C_1$, which just outputs the values of $\bigwedge_{j\in T_i}x_j$   for each $i$.
 
\paragraph{Construction of $C_{i+1}$.} For the inductive case, we proceed as follows. We assume that $i$ is odd (the case that $i$ is even is similar). So by the inductive hypothesis, we know that \expeqref{condition}{eq:oddi} holds and hence that $C_{i}(x) \in I_0(\gamma_i)$ or $I_1(\gamma_i)$ depending on whether $x\in N_\varepsilon$ or $Y_\varepsilon$. Let $\gamma := \gamma_i$. Let the output gates of $C_i$ be $g_1,\ldots,g_M$.

We choose $T_1,\ldots,T_M\subseteq [M]$ randomly as in the statement of \expref{Lemma}{lem:technicalORAND} with $s = A$. Note that the chosen parameters satisfy all the hypotheses of \expref{Lemma}{lem:technicalORAND}. Further we also have
\[
  s\gamma \leq A\cdot A^i\gamma_0 \leq A^{d-2}\cdot \frac{\varepsilon}{\sqrt{n}}\leq \varepsilon_0.
\]

The random circuit $C'$ is defined to be the circuit obtained by adding $M$ OR gates to $C_i$ such that the $j$th OR gate computes $\bigvee_{k\in T_j}g_k$. Let $b_j^x$ be the output of the $j$th OR gate on $C_i(x)$. 

By \expref{Lemma}{lem:technicalORAND}, we have
\begin{align}
x\in N_\varepsilon &\Rightarrow \prob{S}{b_j^x = 0} \geq \exp(-A)\cdot (1+A\gamma\exp(-A\gamma)), \text{ and }\notag \\
x\in Y_\varepsilon &\Rightarrow \prob{S}{b_j^x = 0} \leq \exp(-A)\cdot(1-A\gamma\exp(-A\gamma)).\label{eq:Ci.1}
\end{align}

Let $\delta = {1}/{n^3}$. Note that
\[
  A\gamma \in \left[\frac{1}{\sqrt{n}},\frac{1}{n^{1/2(d-1)}}\right]
\]
and hence for large enough $n$,
\[
  \frac{2\delta}{A\gamma\exp(-A\gamma)}\leq \varepsilon_0.
\]

Assume $x\in N_\varepsilon$. In this case, the Chernoff bound implies that the probability that
\[
  \sum_{j\in [M]}b_j^x\leq M\exp(-A)\cdot (1+A\gamma\exp(-A\gamma))(1-\delta)
\]
is at most $\exp(-\Omega(\delta^2M/\eee^A))\leq \exp(-n)$. When this
event does not occur, we have
\begin{align}
\frac{\sum_i b_i^x}{M}&\geq \frac{1}{\eee^A}(1+A\gamma\exp(-A\gamma))(1-\delta)\notag\\
&\geq \frac{1}{\eee^A}(1+A\gamma\exp(-A\gamma) -2\delta) = \frac{1}{\eee^A}\left(1+A\gamma\exp(-A\gamma)\left(1-\frac{2\delta}{A\gamma\exp(-A\gamma)}\right)\right)\notag\\
&\geq \frac{1}{\eee^A}\left(1+A\gamma\exp(-A\gamma)\cdot \exp\left(-\frac{4\delta}{A\gamma\exp(-A\gamma)}\right)\right)\notag\\
&\geq \frac{1}{\eee^A}(1+A\gamma\exp(-A\gamma)\cdot \exp(-A\gamma))\notag\\
&\geq \frac{1}{\eee^A}(1+A\gamma\exp(-2A\gamma))\geq \frac{1}{\eee^A}(1+\gamma_{i+1}).\label{eq:Ci.2}
\end{align}

We have used above that for large enough $n$,
\[
  \frac{2\delta}{A\gamma\exp(-A\gamma)}\leq \varepsilon_0
\]
and hence
\[
  1-\frac{2\delta}{A\gamma\exp(-A\gamma)}\geq \exp\left(\frac{-4\delta}{A\gamma\exp(-A\gamma)}\right).
\]

Similarly when $x\in Y_\varepsilon$, the Chernoff bound tells us that the probability that
\[
  \sum_{j\in [M]}b_j^x\geq M\exp(-A)\cdot (1-A\gamma\exp(-A\gamma))(1+\delta)
\]
is at most $\exp(-n)$. In this case, we get
\begin{align}
\frac{\sum_i b_i^x}{M}&\leq \frac{1}{\eee^A}(1-A\gamma\exp(-A\gamma))(1+\delta)\notag\\
&\leq \frac{1}{\eee^A}(1-A\gamma\exp(-A\gamma) +\delta) = \frac{1}{\eee^A}\left(1-A\gamma\exp(-A\gamma)\left(1-\frac{\delta}{A\gamma\exp(-A\gamma)}\right)\right)\notag\\
&\leq \frac{1}{\eee^A}\left(1-A\gamma\exp(-A\gamma)\cdot\exp\left(-\frac{2\delta}{A\gamma\exp(-A\gamma)}\right)\right)\notag\\
&\leq \frac{1}{\eee^A}(1-A\gamma\exp(-A\gamma)\cdot \exp(-A\gamma))\notag\\
&= \frac{1}{\eee^A}(1-A\gamma\exp(-2A\gamma))\geq \frac{1}{\eee^A}(1-\gamma_{i+1}).\label{eq:Ci.3}
\end{align}

By a union bound, we can fix $T_1,\ldots,T_M$ so that \expeqref{inequality}{eq:Ci.2} and \expeqref{inequality}{eq:Ci.3} are true for all $x\in N_\varepsilon$ and $x\in Y_\varepsilon$, %
respectively. This gives us the circuit $C_{i+1}$ which satisfies all the required properties.

\paragraph{The top two levels of the circuit.} At the end of the above procedure we have a circuit $C_{d-2}$ of depth $d-2$ and at most $(d-2)M$ gates that satisfies one of \expeqref{condition}{eq:oddi} or \expeqref{condition}{eq:eveni} depending on whether $d-2$ is odd or even,  %
respectively. We assume that $d-2$ is even (the other case is similar). 

Define $\gamma := \gamma_{d-2}$. Recall from \expeqref{inequality}{eq:gamma-i} that $\gamma \geq A^{d-2}\gamma_0\exp(-3A^{d-2}\gamma_0)\geq A^{d-2}\gamma_0/2$.

Let
\[
  M' = \left\lceil \exp\left(\frac{10A\log(1/\varepsilon)}{\varepsilon}+10A\right)\right\rceil.
\]
We choose $M'$ many subsets $T_1,\ldots,T_{M'}\subseteq [M]$ i.i.d.\
so that each $T_j$ is picked as in \expref{Lemma}{lem:technicalORAND}
with $s = 10A\log(1/\varepsilon)/\varepsilon$. Note that

\[
s\gamma \geq s\frac{A^{d-2}\gamma_0}{2}= \frac{10A\log(1/\varepsilon)}{\varepsilon}\cdot \frac{A^{d-2}}{2}\cdot \frac{\varepsilon}{\sqrt{n}}\geq 4\log(1/\varepsilon).
\]
where the final inequality uses the fact that $A^{d-1}= (n^{1/2(d-1)}-\alpha)^{d-1} \geq \sqrt{n}\cdot(1-o(1))$ as long as $2d\leq \log n/(\log \log n + c)$ for a large enough absolute constant $c > 0$.

Say $g_1,\ldots,g_M$ are the output gates of $C_{d-2}$. We define the random circuit $C'$ (with $n$ inputs and $M'$ outputs) to be the circuit obtained by adding $M'$ AND gates such that the $j$th AND gate computes $\bigwedge_{k\in T_j}g_k$. Let $b_j^x$ be the output of the $j$th AND gate on $C_{d-2}(x)$.

By \expref{Lemma}{lem:technicalORAND}, we have
\begin{align}
x\in N_\varepsilon &\Rightarrow \prob{S}{b_j^x = 1} \leq \exp(-s)\cdot \exp(-s\gamma)\leq \varepsilon^2\cdot\exp(-s), \text{ and } \notag \\
x\in Y_\varepsilon &\Rightarrow \prob{S}{b_j^x = 1} \geq \exp(-s)\cdot\exp(s\gamma/2)\geq \frac{\exp(-s)}{\varepsilon^2}. \label{eq:Ctop.1}
\end{align}

Say $x\in N_\varepsilon$. By a Chernoff bound, the probability that $\sum_{j}b_j^x \geq 2\varepsilon^2 M'\exp(-s)$ is at most
\[
  \exp(-\Omega(\varepsilon^2 M'\exp(-s))) \leq \exp(-\Omega(\varepsilon^2 \eee^{10A}))\leq \exp(-n).
\]
Similarly, when $x\in Y_\varepsilon$, the probability that
\[
  \sum_{j}b_j^x \leq\frac{M'\exp(-s)}{2\varepsilon^2}
\]
is also bounded by $\exp(-n)$. By a union bound, we can fix a
$T_1,\ldots,T_{M'}$ to get a circuit $C_{d-1}$ such that
\begin{align}
x\in N_\varepsilon &\Rightarrow |C_{d-1}(x)|_1 \leq 2\varepsilon^2 \exp(-s)M', \text{ and } \notag \\
x\in Y_\varepsilon &\Rightarrow |C_{d-1}(x)|_1 \geq \frac{1}{2\varepsilon^2} \exp(-s)M'. \label{eq:Ctop.2}
\end{align}

This gives us the depth-$(d-1)$ circuit $C_{d-1}$. Note that $C_{d-1}$ has $M' + O(dM) = O(M')$ gates.

To get the depth-$d$ circuit, we choose a random subset $T\subseteq [M']$ by sampling exactly $\lceil\exp(s)\rceil$  many elements of $[M']$ with replacement. We construct a random depth-$d$ circuit $C'_d$ by taking the OR of the output gates of $C_{d-1}$ indexed by the subset $T$. 

{}From \expeqref{implication}{eq:Ctop.2} it follows that 
\begin{align*}
x\in N_\varepsilon &\Rightarrow \prob{T}{C'_d(x) = 1}\leq |T|\cdot 2\varepsilon^2\exp(-s) \leq 4\varepsilon^2 < \varepsilon, \text{ and }\\
x\in Y_\varepsilon &\Rightarrow \prob{T}{C'_d(x) = 0}\leq \left(1-\frac{\exp(-s)}{2\varepsilon^2}\right)^{\exp(s)} \leq \exp(-1/2\varepsilon^2) < \varepsilon.
\end{align*}

The final inequalities in each case above hold as long as $\varepsilon$ is a small enough constant. 

It follows from the above that there is a choice for $T$ such that $C'_d$ makes an error---\ie, $C'_d(x) = 1$ for $x\in N_\varepsilon$ or $C'_d(x) = 0$ for $x\in Y_\varepsilon$---on at most a $2\varepsilon$ fraction of inputs from  $N_\varepsilon\cup Y_\varepsilon$. We fix such a choice for $T$ and the corresponding circuit $C$. 

We have 
\begin{align*}
\prob{x\in \{0,1\}^n}{C(x) \neq \mathrm{Maj}_n(x)}&\leq \prob{x\in Y_\varepsilon\cup N_\varepsilon}{C(x) \neq \mathrm{Maj}_n(x)} + \prob{x\in\{0,1\}^n}{x\not\in Y_\varepsilon\cup N_\varepsilon}\\
&\leq 2\varepsilon + \prob{x\in\{0,1\}^n}{x\not\in Y_\varepsilon\cup N_\varepsilon}.
\end{align*}

Finally by Stirling's approximation we get
\[
\prob{x\in\{0,1\}^n}{x\not\in Y_\varepsilon\cup N_\varepsilon} = \frac{1}{2^n}\sum_{m\in [\frac{n}{2}-\varepsilon\sqrt{n}, \frac{n}{2}+\varepsilon\sqrt{n}]}\binom{n}{m} \leq \frac{1}{2^n}\sum_{m\in [\frac{n}{2}-\varepsilon\sqrt{n}, \frac{n}{2}+\varepsilon\sqrt{n}]}\binom{n}{n/2} \leq 2\varepsilon.
\]

Hence we see that the circuit $C$ computes a $(4\varepsilon,n)$-Approximate Majority, which proves \expref{Theorem}{thm:amano-ckt}.

The circuit has depth $d$ and size $O(M') = \exp(O(n^{1/2(d-1)}\log(1/\varepsilon)/\varepsilon))$.
\end{proof}




\section{Conclusions}
\label{sec:conclusion}

%
%
It is straightforward to extend our main results  %
to $\AC^0[\mathrm{MOD}_p]$ for any fixed prime $p$ to yield a lower bound
of $\exp(\Omega(dn^{1/2(d-1)}/p))$ (the $\Omega(\cdot)$ hides an absolute constant.). To prove this, we construct approximating polynomials as in \expref{Section}{sec:lbd} for size-$s$ depth-$d$ $\AC^0[\mathrm{MOD}_p]$ formulas of degree $O(({1}/{d})p\log s)^{d-1}$. The construction is almost exactly the same,
except that we use $\F_p$-analogs of \expref{Lemma}{lem:razb-OR} (which is
true with the weaker degree bound of
$(p-1)\cdot \lceil \log(1/\varepsilon)\rceil$~\cite{Smolensky87})
and \expref{Lemma}{la:eps} (which holds as stated over all fields~\cite{HS}).
%
The Smolensky--Szegedy  %
degree lower bound for approximating the Majority function
(\expref{Lemma}{lem:razb-maj}) is true over all fields.

The improved construction of approximating polynomials for $\AC^0[\mathrm{MOD}_p]$ formulas also implies another result. Using the fact~\cite{Smolensky87} that any $(1/4)$-approximating polynomial over $\F_p$ ($p$ odd) for the Parity function on $n$ variables must have degree $\Omega(\sqrt{n})$, we see that any polynomial-sized $\AC^0[\mathrm{MOD}_p]$ formula computing the Parity function on $n$ variables must have depth $\Omega(\log n)$ for constant primes $p$. This strengthens a result of Rossman~\cite{Rossman2015average} which gives this statement for $\AC^0$ formulas.


\paragraph{Acknowledgements.}

We thank Rahul Santhanam for valuable discussions, and the organizers of the 2016 Complexity Semester at St.\ Petersburg, where this collaboration began. We are also grateful to the anonymous referees who reviewed this paper for \emph{Theory of Computing} for their corrections and feedback.


%

\nocite{savicky2000number}

%% Theory of Computing
%% File: v015a017.bbl
%% Publication date: December 17, 2019
%% Authors: Ben Rossman and Srikanth Srinivasan
% Define empty bibhead if not already defined
\providecommand{\bibhead}[1]{}
% Define tocrefpdfbookmark if not already defined
\expandafter\ifx\csname pdfbookmark\endcsname\relax%
  \providecommand{\tocrefpdfbookmark}{}
\else\providecommand{\tocrefpdfbookmark}{%
   \hypertarget{tocreferences}{}%
   \pdfbookmark[1]{References}{tocreferences}}%
\fi

\tocrefpdfbookmark
\begin{thebibliography}{10}

\bibitem{Amano-appxmaj}\bibhead{Amano-appxmaj}
{\sc Kazuyuki Amano}: Bounds on the size of small depth circuits for
  approximating majority.
\newblock In {\em Proc. 36th Internat. Colloq. on Automata, Languages and
  Programming (ICALP'09)}, pp. 59--70. Springer, 2009.
\newblock [\epfmtdoi{10.1007/978-3-642-02927-1\_7}]

\bibitem{AroraBarak}\bibhead{AroraBarak}
{\sc Sanjeev Arora and Boaz Barak}: {\em Computational Complexity - {A} Modern
  Approach}.
\newblock Cambridge Univ. Press, 2009.
\newblock [\epfmtdoi{10.1017/CBO9780511804090}]

\bibitem{Beigelpolysurvey}\bibhead{Beigelpolysurvey}
{\sc Richard Beigel}: The polynomial method in circuit complexity.
\newblock In {\em Proc. 8th IEEE Conf. on Structure in Complexity Theory
  (SCT'93)}, pp. 82--95. IEEE Comp. Soc. Press, 1993.
\newblock [\epfmtdoi{10.1109/SCT.1993.336538}]

\bibitem{BlaisTan}\bibhead{BlaisTan}
{\sc Eric Blais and Li{-}Yang Tan}: Approximating {B}oolean functions with
  depth-2 circuits.
\newblock {\em SIAM J. Comput.}, 44(6):1583--1600, 2015.
\newblock Preliminary version in
  \href{http://doi.org/10.1109/CCC.2013.17}{CCC'13}.
\newblock [\epfmtdoi{10.1137/14097402X}]

\bibitem{HS}\bibhead{HS}
{\sc Prahladh Harsha and Srikanth Srinivasan}: On polynomial approximations to
  {$\mathrm{AC}^0$}.
\newblock {\em Random Structures Algorithms}, 54(2):289--303, 2019.
\newblock Preliminary version in
  \href{http://doi.org/10.4230/LIPIcs.APPROX-RANDOM.2016.32}{RANDOM'16}.
\newblock [\epfmtdoi{10.1002/rsa.20786}, \epfmt{arxiv}{1604.08121}]

\bibitem{Hastad}\bibhead{Hastad}
{\sc Johan H{\aa}stad}: Almost optimal lower bounds for small depth circuits.
\newblock In {\em Proc. 18th STOC}, pp. 6--20. ACM Press, 1986.
\newblock [\epfmtdoi{10.1145/12130.12132}]

\bibitem{hoory2006monotone}\bibhead{hoory2006monotone}
{\sc Shlomo Hoory, Avner Magen, and Toniann Pitassi}: Monotone circuits for the
  majority function.
\newblock In {\em Proc. Internat. Workshop on Approximation, Randomization, and
  Combinatorial Optimization. Algorithms and Techniques (RANDOM'06)}, pp.
  410--425. Springer, 2006.
\newblock [\epfmtdoi{10.1007/11830924\_38}]

\bibitem{Jukna}\bibhead{Jukna}
{\sc Stasys Jukna}: {\em Boolean Function Complexity: Advances and Frontiers}.
\newblock Springer, 2012.
\newblock [\epfmtdoi{10.1007/978-3-642-24508-4}]

\bibitem{karchmer1990monotone}\bibhead{karchmer1990monotone}
{\sc Mauricio Karchmer and Avi Wigderson}: Monotone circuits for connectivity
  require super-logarithmic depth.
\newblock {\em SIAM J. Discrete Math.}, 3(2):255--265, 1990.
\newblock Preliminary version in
  \href{http://doi.org/10.1145/62212.62265}{STOC'88}.
\newblock [\epfmtdoi{10.1137/0403021}]

\bibitem{KPPY84majubd}\bibhead{KPPY84majubd}
{\sc Maria~M. Klawe, Wolfgang~J. Paul, Nicholas Pippenger, and Mihalis
  Yannakakis}: On monotone formulae with restricted depth.
\newblock In {\em Proc. 16th STOC}, pp. 480--487. ACM Press, 1984.
\newblock [\epfmtdoi{10.1145/800057.808717}]

\bibitem{kopparty2012certifying}\bibhead{kopparty2012certifying}
{\sc Swastik Kopparty and Srikanth Srinivasan}: Certifying polynomials for
  $\mathrm{AC}^0[\oplus]$ circuits, with applications to lower bounds and
  circuit compression.
\newblock {\em Theory of Computing}, 14(12):1--24, 2018.
\newblock Preliminary version in
  \href{http://drops.dagstuhl.de/opus/volltexte/2012/3846/}{FSTTCS'12}.
\newblock [\epfmtdoi{10.4086/toc.2018.v014a012}]

\bibitem{odonnellwimmer}\bibhead{odonnellwimmer}
{\sc Ryan O'Donnell and Karl Wimmer}: Approximation by {DNF:} {E}xamples and
  counterexamples.
\newblock In {\em Proc. 34th Internat. Colloq. on Automata, Languages and
  Programming (ICALP'07)}, pp. 195--206. Springer, 2007.
\newblock [\epfmtdoi{10.1007/978-3-540-73420-8\_19}]

\bibitem{Razborov}\bibhead{Razborov}
{\sc Alexander~Alexandrovich Razborov}: Lower bounds on the size of bounded
  depth circuits over a complete basis with logical addition.
\newblock {\em Mathematical Notes of the Academy of Sciences of the USSR},
  41(4):333--338, 1987.
\newblock Russian original in \emph{Mat. Zametki} 41(4), 1987, 598--607.
\newblock [\epfmtdoi{10.1007/BF01137685}]

\bibitem{Rossman2015average}\bibhead{Rossman2015average}
{\sc Benjamin Rossman}: The average sensitivity of bounded-depth formulas.
\newblock {\em Comput. Complexity}, 27(2):209--223, 2018.
\newblock Preliminary version in
  \href{http://doi.org/10.1109/FOCS.2015.33}{FOCS'15}.
\newblock [\epfmtdoi{10.1007/s00037-017-0156-0}, \epfmt{arxiv}{1508.07677}]

\bibitem{savicky2000number}\bibhead{savicky2000number}
{\sc Petr Savick\'y and Alan~R. Woods}: The number of {B}oolean functions
  computed by formulas of a given size.
\newblock {\em Random Structures Algorithms}, 13(3--4):349--382, 2000.
\newblock
  [\epfmtdoi{10.1002/(SICI)1098-2418(199810/12)13:3/4<349::AID-RSA9>3.0.CO;2-V}]

\bibitem{Smolensky87}\bibhead{Smolensky87}
{\sc Roman Smolensky}: Algebraic methods in the theory of lower bounds for
  {B}oolean circuit complexity.
\newblock In {\em Proc. 19th STOC}, pp. 77--82. ACM Press, 1987.
\newblock [\epfmtdoi{10.1145/28395.28404}]

\bibitem{Smolensky93}\bibhead{Smolensky93}
{\sc Roman Smolensky}: On representations by low-degree polynomials.
\newblock In {\em Proc. 34th FOCS}, pp. 130--138. IEEE Comp. Soc. Press, 1993.
\newblock [\epfmtdoi{10.1109/SFCS.1993.366874}]

\bibitem{Szegedy1989}\bibhead{Szegedy1989}
{\sc Mario Szegedy}: {\em Algebraic Methods in Lower Bounds for Computational
  Models with Limited Communication}.
\newblock Ph.\,D.\ thesis, University of Chicago, 1989.
\newblock Available at
  \href{https://people.cs.uchicago.edu/~laci/students/szegedy.dir}{advisor's
  home page}.

\bibitem{Valiant-maj}\bibhead{Valiant-maj}
{\sc Leslie~G. Valiant}: Short monotone formulae for the majority function.
\newblock {\em J. Algorithms}, 5(3):363--366, 1984.
\newblock [\epfmtdoi{10.1016/0196-6774(84)90016-6}]

\end{thebibliography}

\end{document}
